%%%%%%%%%%%%%%%%%%%%%%%%%%%%%%%%%%%%%%%%%%%%%%
%% 04_volatility_regimes.tex --- 正文层 (BODY LAYER)
%%
%% 这里只有内容。字体、边距、颜色、定理环境长什么样，全在
%% ../../preamble.tex 里；改版式不用碰这个文件。
%%
%% 这是正文片段，没有 \documentclass 也没有 document 环境：
%% ../../build.sh 单独编译它，all_chapters.tex 把它 \input 进合订本。
%%%%%%%%%%%%%%%%%%%%%%%%%%%%%%%%%%%%%%%%%%%%%%

\chapterhead{04}{Volatility Regimes}{04_volatility_regimes}

Three steps, in order. \textbf{MACD stops working when the tape gets loud}
(\secref{sec:macd-breaks}). \textbf{The option market prices that loudness before any trailing
estimate can measure it} (\secref{sec:option-market}). \textbf{The resulting label is used by
conditioning on it, not by adding it to the signal} (\secref{sec:using-label}).


\section{Why MACD breaks when the tape gets loud}\label{sec:macd-breaks}

\fig{macd-out-of-phase}{Above, one price path with a 12-day and a 26-day EMA over it. Through the
calm first stretch the fast average sits clear above the slow one and the two never cross; through
the shaded stretch, which has no net trend and four times the daily amplitude, they stay tangled and
their difference changes sign ten times. Below, the position the rule holds: one unbroken long, then
a shredded alternation of long and short blocks. Each flip is a round trip, paid at the widest
spreads of the year. Illustrative.}{fig:macd-out-of-phase}

One path, one rule, two regimes. In the calm stretch the fast average pulls clear of the slow one
and stays there: no crossings, one position, held. In the shaded stretch the trend is gone and the
swings are four times wider, so the two legs stay tangled and noise carries one across the other
again and again.

\textbf{Every flip also lands after the turn it is chasing}, since an average only confirms a move
that has already happened---so the rule goes short at the bottom of the spike and long at the top of
the snap-back. Whipsaw is not scatter around the right answer; it is the opposite of it, repeated.

Standardizing the signal (\chapref{\S 2.4}{02_testing_a_signal}) rescales it without moving a
crossing date, and \chapref{\S 3.4}{03_shaping_the_lookback}'s smoother and deadband damp every date
alike because neither knows the tape is violent. Clearing that ceiling means promoting the state of
the market to a \textbf{variable}.


\section{What the option market already knows}\label{sec:option-market}

\subsection{Volatility cannot be computed from the price}\label{sec:vol-not-from-price}

There is no way to read today's volatility off the price series. Every estimate built from past
returns---a rolling standard deviation over a month, an EWMA, any weighting you like---is an average
of what has already happened, so it only climbs as the new regime's returns accumulate inside its
window. Shortening the window trades the lag for noise and buys nothing.

\fig{regime-lag}{Above, the absolute daily return around a news event: quiet either side, with a
burst of three and four percent moves in the sessions after it prints. Below, two answers to that
event in annualized volatility. The implied line jumps on the day the news prints and decays
smoothly back; the trailing 21-day line barely moves on the day and does not peak until session
twenty-one, by which time the burst has been over for a fortnight.
Illustrative.}{fig:regime-lag}

The upper panel is the event and the lower panel is two answers to it. The trailing line does not
peak until three weeks after the burst it is describing. \textbf{The other line is not an estimate
at all}---it is a price, quoted by people who have to take a view on the next thirty days, and it
moves the session the news does. Where it comes from is \secref{sec:option-to-vol}.

\subsection{How an option hands you a volatility}\label{sec:option-to-vol}

An option fixes a price today for a trade that may happen later, and \textbf{the buyer chooses
whether it happens}. Two of them, defined by which side that choice sits on:

\begin{tabularx}{\linewidth}{|>{\raggedright\arraybackslash}p{0.30\linewidth}|X|X|}
\hline
 & \textbf{Call} & \textbf{Put} \\ \hline
The holder may & \textbf{buy} the underlying at $K$ & \textbf{sell} the underlying at $K$ \\ \hline
Worth exercising when & $S_T > K$ & $S_T < K$ \\ \hline
Payoff at expiry & $\max(S_T - K, 0)$ & $\max(K - S_T, 0)$ \\ \hline
Pays off when the underlying & rises & falls \\ \hline
The most the buyer can lose & the premium & the premium \\ \hline
The most the buyer can make & unbounded & $K$, less the premium \\ \hline
Bought by someone who & wants upside without owning it & wants a floor under something they
                                                         own \\ \hline
\end{tabularx}

with $K$ the \textbf{strike}---the price fixed in the contract---and $S_T$ the underlying's price at
expiry.

\fig{option-payoffs}{Value at expiry against the underlying's price at expiry, with the strike $K$
marked in both panels. Left, a call: flat at zero below the strike, rising one-for-one above it, and
a dashed line one premium lower for the same shape net of what was paid. Right, a put: the mirror
image. Both dashed lines are flat on their far side---the loss is capped at the premium.
Illustrative.}{fig:option-payoffs}

Both payoffs are \textbf{kinked}: flat on one side of the strike, one-for-one on the other. That
kink is the whole instrument---it truncates the loss and leaves the gain open---and it is why an
option is worth more when the underlying might travel further. Under Black--Scholes a call is worth
\[
  C = S \Phi(d_1) - K e^{-R \tau} \Phi(d_2), \qquad
  d_1 = \frac{\ln \left( S/K \right) + \left( R + \sigma^2/2 \right) \tau}{\sigma \tau^{1/2}},
  \qquad
  d_2 = d_1 - \sigma \tau^{1/2}
\]
where $S$ is the spot price, $K$ the strike, $\tau$ the time to expiry in years, $R$ the risk-free
rate, $\Phi$ the standard normal cumulative distribution function, and $\sigma$ the volatility.
Everything but $\sigma$ is observable, which is what makes the equation worth inverting.

\begin{definition}[Implied volatility]\label{def:implied-vol}
The value of $\sigma$ that makes the model reproduce the option's traded price---a forecast for the
life of the contract, read out of a price rather than estimated from history.
\end{definition}

\begin{claim}\label{clm:inversion}
The inversion is well defined: one price, one implied volatility.
\end{claim}

\begin{proof}
The kink makes both payoffs convex in $S_T$. Raising $\sigma$ spreads the distribution of $S_T$
while leaving its forward mean where it was, and the expectation of a convex function rises under a
mean-preserving spread. The price is therefore strictly monotone in $\sigma$, and a monotone map is
invertible.
\end{proof}

There is no closed form for the inverse, so $\sigma$ is solved numerically---bisection needs about
twenty steps to reach four decimals, and monotonicity is what guarantees it converges. At the money
the whole thing collapses to something you can do in your head:
\[
  C \approx 0.4\, S \sigma \tau^{1/2}
  \qquad \Longrightarrow \qquad
  \sigma \approx \frac{C}{0.4\, S \tau^{1/2}}
\]

\begin{example}\label{ex:implied-vol}
A 30-day at-the-money call trades at 2 percent of spot. Quote the price as a fraction of spot and
$S$ cancels, leaving a denominator of $0.4 \tau^{1/2} \approx 0.115$ at $\tau = 30/365$, so
$\sigma \approx 0.02/0.115 \approx 0.17$---the market is charging for a \textbf{17 percent
annualized} move over the next month. Read the same contract a day later at 3 percent of spot and
the forecast is 26 percent: the number moves with the quote, which is the entire point.
\end{example}

\begin{note}[One contract is not the market]
That 17 percent is the volatility charged \textbf{at that strike}. Quotes on other strikes imply
different numbers---the smile---so a single contract's answer is a property of one contract.
\secref{sec:vix} reads the whole strip at once to get rid of that dependence.
\end{note}

\subsection{VIX}\label{sec:vix}

\begin{definition}[VIX]\label{def:vix}
A constant-30-day implied volatility for the S\&P 500, built from the whole strip of
out-of-the-money options rather than from any one of them. For a single expiry the strip gives the
implied \textbf{variance} directly, with no model to invert:
\[
  \sigma^2_\tau = \frac{2 e^{R \tau}}{\tau} \sum_i \frac{\Delta K_i}{K_i^2} Q(K_i)
                  - \frac{1}{\tau} \left( \frac{F}{K_0} - 1 \right)^2
\]
where $Q(K_i)$ is the mid price of the out-of-the-money option struck at $K_i$, $\Delta K_i$ the gap
between neighbouring strikes, $F$ the forward level implied by put-call parity, and $K_0$ the first
strike below $F$. Every term is a traded price, so the sum is the market's own variance rather than
one contract's.
\end{definition}

Two steps turn that into the published number. \textbf{Take the two expiries with $\tau$ between 23
and 37 days}---listed contracts do not fall on a 30-day schedule, and the band is what guarantees
one on each side of the target---and \textbf{interpolate their variances linearly in $\tau$} onto
exactly 30 days. Then annualize, take the root, and quote it in percentage points:
\[
  \mathrm{VIX}_t = 100 \left( \sigma^2_{30,t} \right)^{1/2}
\]

\begin{note}[Decide whether your formula wants the variance or the volatility]
The published level is already a volatility---a VIX of 20 means an annualized 20 percent---because
the strip computes a variance and the last step takes its root. Variance is what is additive, across
time and across independent sources of risk; volatility is not. Anything that rescales a horizon or
adds contributions needs $\left( \mathrm{VIX}_t / 100 \right)^2$; anything compared against a return
in the same units needs the level. The error is invisible when made: a missing square root still
leaves a plausible positive number.
\end{note}

\begin{note}[What the forecast is worth is a separate question]
Implied volatility is what the market charges, not what will happen---insurance is sold above its
expected cost, and that gap is a strategy in its own right. Here the index is only a \textbf{state
label}, for which a live consensus is enough.
\end{note}

\subsection{One index for equities, another for rates}\label{sec:two-indices}

\begin{claim}\label{clm:one-index}
One equity index's implied volatility labels most sleeves, even though it is computed from one
market.
\end{claim}

The reason is what a panic does to correlations. In calm markets the sleeves move mostly apart; in a
large enough shock the same levered participants sell all of them at once to raise cash and the
cross-asset correlations run toward one. The states in which the label matters most are exactly the
states in which one market's fear is every market's fear---an empirical regularity rather than a
theorem, and better at the extremes than in the middle.

\begin{tabularx}{\linewidth}{|>{\raggedright\arraybackslash}p{0.30\linewidth}|>{\raggedright\arraybackslash}p{0.14\linewidth}|X|}
\hline
\textbf{Sleeve} & \textbf{Regime index} & \textbf{What its spikes are about} \\ \hline
Equities, credit, most commodities & \textbf{VIX}
  & growth and earnings shocks, and forced deleveraging \\ \hline
Rates, and anything priced off the curve & \textbf{MOVE}
  & policy surprises, inflation prints, auction stress \\ \hline
\end{tabularx}

\begin{note}[Rates are the genuine exception]
Treasury volatility is close to unrelated to equity volatility, because the two are frightened by
different things. A rate sleeve labelled by VIX is labelled wrongly, and \textbf{MOVE} is the same
construction run on Treasury options. Neither index is in \texttt{CTA\_data/}
(\chapref{100 \midot{} The Dataset}{100_dataset}), so both have to be fetched before
\secref{sec:using-label} can be run.
\end{note}

\subsection{From a level to a label}\label{sec:level-to-label}

\begin{definition}[Volatility regime]\label{def:vol-regime}
A stretch of dates over which the amplitude of returns is roughly constant and materially different
from the stretch on either side; a \textbf{regime shift} is the boundary between two of them.
\end{definition}

\secref{sec:using-label} needs a small set of such states, named, because a continuous variable
cannot be a row of a table. Write $v_t$ for the index level on date $t$; its distribution is the
obstacle.

\fig{vix-bucketing}{The same illustrative sample of volatility index levels, twice. Left, against
the raw level: the mass piles into a narrow band and trails off into a long near-empty tail, half
the sample sitting inside six points. Right, against the natural logarithm of the same values: close
to symmetric, and the two integer cuts at 3 and 4---about 20 and about 55 on the level---divide it
into calm, stressed and crisis.}{fig:vix-bucketing}

The level spends most of its life in a band a few points wide and occasionally reaches four times
that, so equal-width bins put nearly everything in the first one---the opposite of what
\chapref{\S 2.3.2.2}{02_testing_a_signal}'s bucket test needs.

\begin{definition}[Log-ceiling bucket]\label{def:log-ceiling}
Take the natural log of the level and round up to the next integer,
$g_t = \left\lceil \ln v_t \right\rceil$, where $g_t$ is the regime label on that date. Because
$\ln$ compresses the sparse upper tail and stretches the crowded lower band, the integer cuts land
at multiplicatively spaced levels:

\begin{center}
\begin{tabular}{|c|l|l|}
\hline
\textbf{Label $g_t$} & \textbf{Level range} & \textbf{Reads as} \\ \hline
3 & $v_t \in (7.4, 20.1]$   & calm---the ordinary tape \\ \hline
4 & $v_t \in (20.1, 54.6]$  & stressed \\ \hline
5 & $v_t \in (54.6, 148.4]$ & crisis \\ \hline
\end{tabular}
\end{center}
\end{definition}

The boundaries are $e^3$ and $e^4$, close to where a practitioner would have drawn them by hand.
That is the construction's virtue: \textbf{the cuts use no data, so a date's label means the same
thing in 2015 and in 2020.} Writing $g_t = \left\lceil c \ln v_t \right\rceil$ puts the boundaries
at $e^{k/c}$ and recovers as many buckets as wanted; $c$ is a hyperparameter like any other, and
choosing it by which value produces the best-looking grid is what
\chapref{5 \midot{} Modelling}{05_modelling} is about.

The alternative is \chapref{\S 2.3.2.1}{02_testing_a_signal}'s own machinery applied to the index:
rank the dates by $v_t$ and cut at the quintiles.

\begin{tabularx}{\linewidth}{|>{\raggedright\arraybackslash}p{0.24\linewidth}|X|X|}
\hline
 & \textbf{Log-ceiling} & \textbf{Rolling quantile} \\ \hline
\textbf{Boundaries}            & fixed, the same in every year
                               & move with the window \\ \hline
\textbf{Group sizes}           & very unequal---crisis may be 3 percent of the sample
                               & equal by construction \\ \hline
\textbf{A label means}         & an absolute level of market fear
                               & a rank against the recent past \\ \hline
\textbf{Comparable across time} & yes
                               & no---the same $v_t$ can be G5 one year and G3 the next \\ \hline
\textbf{Look-ahead risk}       & none, the cuts use no data
                               & real, if the window is not trailing \\ \hline
\end{tabularx}

\begin{note}[The ranking must be rolling, never full-sample]
A quintile computed over the whole history uses 2020's levels to label 2015's dates, which is
look-ahead bias of the plainest kind (\chapref{\S 2.6}{02_testing_a_signal})---and it is easy to
commit here, because the index feels like context rather than like data.
\end{note}


\section{Using the label}\label{sec:using-label}

Each date now carries three numbers: a signal value, a forward return, and a regime label $g_t$. Two
different questions can be asked of that---whether the edge \textbf{differs} by regime
(\secref{sec:conditioning}), and whether the regime is \textbf{new information} the signals do not
already carry (\secref{sec:residualizing}).

\subsection{Conditioning --- the two-way sort}\label{sec:conditioning}

Put the signal buckets across and the regime labels up, and report the mean forward return in every
cell. It is \chapref{\S 2.3.2}{02_testing_a_signal}'s bar plot run once per regime---equivalently a
\textbf{conditional model}: sort on the state first, then on the signal within it.

\fig{regime-signal-grid}{Mean forward return in basis points, signal buckets G1 to G5 across and
volatility regime up: calm, stressed, crisis. The calm row, holding three-quarters of the sample,
rises monotonically; the stressed row rises about half as steeply; the crisis row, 3 percent of the
sample, has no ordering at all and its shading is flat across the row.
Illustrative.}{fig:regime-signal-grid}

Read it by rows, and the reading is the finding. A staircase in the calm row that flattens through
the stressed row and disappears in the crisis row is \secref{sec:macd-breaks}'s picture measured on
your own data: same rule, same parameters, less and less to show for them. The one-dimensional plot
of \chapref{2}{02_testing_a_signal} is the sample-weighted average of the three rows---it reports
the calm row lightly diluted and never mentions that the other two exist.

\begin{note}[Keep the regime axis coarse, and print the counts]
The cells partition the same sample, so a one-way bar's $N/5$ dates become $N/15$ across three rows
and every error bar (\chapref{\S 2.3.2.3}{02_testing_a_signal}) widens by
$3^{1/2} \approx 1.7$---at five rows, $5^{1/2} \approx 2.2$. Worse, the row you most want is the
thinnest: crisis cells hold roughly $N/167$ dates, error bars about 5.8 times wider, so \textbf{a
flat crisis row is equally consistent with an edge nobody has enough crisis days to measure.} Report
the count beside every mean, and treat a row you cannot measure as unknown rather than as zero.
\end{note}

\subsection{Residualizing --- is any of it new?}\label{sec:residualizing}

Regress the candidate on the incumbents and keep what is left:
\[
  x_t = \beta_0 + \sum_j \beta_j x_{j,t} + \epsilon_t
\]
where $x_t$ is the candidate signal on date $t$, $x_{j,t}$ the signals already in the book
---momentum and MACD, here---$\beta_j$ their fitted coefficients, and $\epsilon_t$ the residual.
Then run \chapref{\S 2.3.2}{02_testing_a_signal}'s bar plot on $\epsilon_t$ in place of $x_t$.

\begin{claim}\label{clm:residual}
If the staircase survives on $\epsilon_t$, the candidate carries an edge the existing signals do not
already own.
\end{claim}

\begin{proof}
Least squares sets the derivative of the summed squared error to zero, giving one normal equation per
regressor---$\sum_t \epsilon_t x_{j,t} = 0$ for every $j$, and $\sum_t \epsilon_t = 0$ from the
intercept. Zero sum of products on a centred variable is zero sample covariance, so
$\rho(\epsilon, x_j) = 0$ for every incumbent. By \chapref{\S 2.1.2}{02_testing_a_signal} a
portfolio's expected return runs entirely through the correlation between what sets the weights and
the forward return, so weights built on $\epsilon_t$ earn whatever they earn somewhere the book is
not standing.
\end{proof}

\begin{note}[Orthogonal is not independent]
Zero \emph{linear} correlation is all least squares buys. A residual that is large exactly when MACD
is large \textbf{in absolute value} is a function of MACD the regression cannot see---and that is
not a corner case here, since the candidate is a variable about $\sigma$ and
$\left| \mathrm{MACD} \right|$ widens with $\sigma$ too. Add the transform you suspect as an extra
regressor and residualize against that too.
\end{note}

\begin{note}[With fifty signals, regress against the model's output]
Fitting a candidate against fifty regressors on overlapping daily data fits noise long before it
exhausts the information, and fifty pairwise regressions give answers that do not compose. The book
already produces one number per asset per date---its combined prediction $p_t$. Use that single
series as the regressor: one residual, and the question it answers is the one actually being asked.
\end{note}

\subsection{Where the answer is spent}\label{sec:answer-spent}

\textbf{Often not in the return model at all.} The natural home for \secref{sec:conditioning}'s grid
is the \textbf{optimizer}---a penalty that rises with the regime, an exposure cap that binds in the
crisis row, a volatility target that scales the whole book---which is why the grid earns its place
even when \secref{sec:residualizing}'s residual comes back flat. Answering \emph{yes} to
\secref{sec:conditioning} and \emph{no} to \secref{sec:residualizing} is a common and perfectly good
result, and performance evaluation is where it is spent.

\textbf{Add one dimension per experiment, and record what it bought}---what the step adds alone,
whether it overlaps what is already there (\secref{sec:residualizing}'s residual), and which
direction it is useful in, since a variable can predict return poorly and volatility well. The
record has to pass interpretability in an operational sense: move an input and you should be able to
say in advance, roughly, what the output does. Measuring 8 bp where you expected 10 is fine;
\emph{``it moved and I do not know by how much''} is not, and it is what adding several things at
once guarantees.

\begin{note}[The scorecard is regime-dependent too]
Every chart above reports mean forward return with an error bar, for the reasons
\chapref{\S 2.5}{02_testing_a_signal} sets out. A Sharpe ratio has the tape's volatility in its own
denominator, so the same number means different things in different rows: 0.8 earned in a year when
the index itself ran at 0.4 is a far harder result than 0.9 earned in a year when buying the index
and doing nothing paid 0.8. Collapsing the regime axis into one ratio destroys exactly the context
needed to read the ratio.
\end{note}


\section*{Background}
\addcontentsline{toc}{section}{Background}

\subsection*{Why is an option the instrument a volatility forecast comes from?}

\begin{itemize}
  \item \textbf{Its value depends on the range, not the direction.} The flat half of the payoff
        truncates the losses, so a wider distribution of outcomes is worth strictly more---which
        makes an option the most volatility-sensitive thing quoted, and
        \secref{sec:option-to-vol}'s inversion stable enough to be worth doing.
  \item \textbf{Out-of-the-money contracts are close to pure volatility bets.} One struck far above
        the price pays only if the underlying travels an unusual distance, so its price is almost
        entirely a statement about the tails---which is why VIX reads the whole strip, not the
        at-the-money contract.
  \item \textbf{The forecast has a fixed horizon attached.} Every contract has an expiry, so an
        implied volatility is always \emph{volatility over the next $\tau$ days}---which is what
        allows two of them to be interpolated onto a constant 30-day horizon.
\end{itemize}


\section*{Appendix \midot{} Notation}
\addcontentsline{toc}{section}{Appendix \midot{} Notation}

Throughout, $t$ is the date. The regime index is a property of the market rather than of one asset,
so it carries no asset subscript $s$.

\begin{xltabular}{\linewidth}{|>{\raggedright\arraybackslash}p{0.22\linewidth}|X|>{\raggedright\arraybackslash}p{0.16\linewidth}|}
\hline
\textbf{Symbol} & \textbf{Means} & \textbf{First used} \\ \hline
\endhead
$\sigma$
  & the volatility of the underlying---the amplitude of its move, without regard to sign
  & \secref{sec:option-to-vol} \\ \hline
$C$, $S$, $K$, $\tau$, $R$
  & a call's price, the spot price, the strike, the time to expiry in years, and the risk-free rate
  & \secref{sec:option-to-vol} \\ \hline
$\Phi$, $d_1$, $d_2$
  & the standard normal cumulative distribution function and the two arguments Black--Scholes feeds
    it
  & \secref{sec:option-to-vol} \\ \hline
$S_T$
  & the underlying's price at expiry
  & \secref{sec:option-to-vol} \\ \hline
$Q(K_i)$, $\Delta K_i$, $F$, $K_0$
  & one out-of-the-money option's mid price, the gap between neighbouring strikes, the forward
    level, and the first strike below it
  & \secref{sec:vix} \\ \hline
$\sigma^2_\tau$, $\sigma^2_{30,t}$
  & the implied variance the strip gives for one expiry, and the same interpolated onto 30 days
  & \secref{sec:vix} \\ \hline
$v_t$, $g_t$, $c$
  & the index level on that date, the regime label cut from it, and the multiplier setting how fine
    the cuts are
  & \secref{sec:level-to-label} \\ \hline
$x_t$, $x_{j,t}$, $\beta_j$, $\epsilon_t$
  & the candidate signal, the signals already in the book, their fitted coefficients, and the
    residual
  & \secref{sec:residualizing} \\ \hline
$p_t$, $N$
  & the model's combined prediction, and the number of dates in the sample
  & \secref{sec:using-label} \\ \hline
\end{xltabular}

\begin{note}[Collisions to watch]
$\sigma$ is a market-wide volatility here where \chapref{\S 2.4}{02_testing_a_signal}'s
$\sigma_{s,t}$ is one asset's trailing estimate, and $R$ is the risk-free rate discounting an option
where \chapref{\S 2.5}{02_testing_a_signal}'s $r_f$ is the same rate subtracted from a portfolio
return. Uppercase $S$ is a spot price and lowercase $s$ is
\chapref{2}{02_testing_a_signal}'s asset index; $\Phi$ is a normal CDF and $N$ a sample size, which
is why the CDF is written $\Phi$ rather than $N$. Lowercase $g_t$ is a regime label on the vertical
axis while \chapref{\S 2.3.2}{02_testing_a_signal}'s uppercase G1 \ldots G5 are signal buckets on
the horizontal one, and \secref{sec:conditioning}'s grid has both at once. $\epsilon_t$ here is a
signal residual after regressing on other \textbf{signals};
\chapref{\S 2.2}{02_testing_a_signal}'s $\epsilon$ is the part of the forward \textbf{return} no
signal reaches.
\end{note}
